% NOETHER PAPER 4 — KOREAN TRANSLATION-PRODUCER DRAFT — UNCHECKED
% Unit: T03-U16; current-authority whole lines 3865--3886
% Source slice: 2,161 LF UTF-8 bytes; SHA-256 7860EA7EA30526C7F93C223463B78E6054C9A28F7C7F926A9E4C37CAD9A2AC19
% Pointer: NOETH-DE-AUTH-v004-20260804; SHA-256 A1C62FDACAA34DFC1B806DC18258F2E732539F7AE9D85AA4BD9E1067B8749D9F
% This is producer translation only: no source/Korean/formula review, compilation, or rendering.

정리의 첫째 부분은 (20.)과 (21.)을 유도할 때의 고찰에서 나오며, 둘째 부분은 다음과 같은 두 기호열 사이의 항등식이 존재할 수 없다는 데서 나온다.
\[
\pair{q_\rho A^\sigma}{B_{\rho+\sigma-\tau}p_\tau}
=\eps_1\pair{q_\rho B^{n-\rho-\sigma+\tau}}{A_{n-\sigma}p_\tau}
+\eps_2\pair{q_\rho q_{n-\tau}}{B_{\rho+\sigma-\tau}A_{n-\sigma}},
\]
여기서 $\rho\leqq\tau\leqq\sigma$이고 어느 열도 무게 $1\cdot(n-1)$을 갖지 않는다. ($\rho,\tau,\sigma$에 관한 다른 가정은 열의 표기를 서로 바꾸면 이 경우로 환원된다.) 예를 들어 이 불가능성은 (8.)을 전개하되 $q_\rho=lM^{\rho-1}$, $q_{n-\tau}=lN^{n-\tau-1}$로 특수화하여 얻는다. 이때 $l_1=1$, $M^{2,3,\ldots,\rho}=1$, $A^{\rho+1,\ldots,\rho+\sigma}=1$로 놓고 열 $l,M,A$의 나머지 모든 원소를 0으로 놓는 한편, $N$과 $B$는 임의로 둔다. 이제 임의의 열 사이의 항등식은 하나의 열 $p_\tau$를 $y^{(1)}y^{(2)}\cdots y^{(\tau)}$로 풀어 놓음으로써 (20.)으로부터 합성되는 항등식으로 바뀐다. 그러므로 (20.)을 기호열 두 개만을 가진 항등식들로부터 생성할 수 있음을 보이기만 하면 정리의 둘째 부분이 증명된다. 실제로
\begin{equation}
\pair{A^{\rho_1}A^{\rho_2}}{xp_{\rho-1}}
=\eps\pair{A^{\rho_1}q_{n-\rho+1}}{A_{n-\rho_2}x}
+\pair{A^{\rho_1}}{xp_{\rho_1-1}},
\quad \text{여기서 }p_{\rho_1-1}\sim A^{\rho_2}q_{n-\rho+1},
\tag{20a}
\end{equation}
에서 $A^{\rho_1}=B^{\sigma_1}B^{\sigma_2}$로 특수화한다. 즉,
\[
\pair{B^{\sigma_1}B^{\sigma_2}}{xp_{\rho_1-1}}
=\eps\pair{B^{\sigma_1}q_{n-\rho_1+1}}{B_{n-\sigma_2}x}
+\pair{B^{\sigma_1}}{xp_{\sigma_1-1}},
\quad \text{여기서 } p_{\sigma_1-1}\sim B^{\sigma_2}q_{n-\rho_1+1},
\]
로 특수화하면 세 기호열에 관한 하나의 항등식 (20.)을 얻는다. 그리고 $B^{\sigma_1}=C^{\tau_1}C^{\tau_2}$ 등으로 계속 특수화하면 일반 항등식 (20.)을 얻는다. 따라서 두 기호열과 두 임의의 변수열 사이의 항등식에는 명시적 표현이 존재하지 않는다는 사실로부터, 변수열들 가운데 열 $x$나 $u$가 하나도 없다면 기호열이 임의로 많이 있는 경우에도 명시적 표현이 존재하지 않는다는 결론이 따른다.
